\documentclass[10pt]{amsart}
\usepackage{calc,amssymb,amsmath,ulem,stmaryrd,fullpage}
\usepackage{enumerate}
\usepackage{alltt}
\usepackage{multicol}
\RequirePackage[dvipsnames,usenames]{color}
\let\mffont=\sf
\normalem
%\input{kmacros3.sty}
\input{xy}
\xyoption{all}
\usepackage{hyperref}
\usepackage{graphicx}
\usepackage[all,cmtip]{xy}
\usepackage{verbatim}
\usepackage[left=0.95in,top=0.95in,right=0.95in,bottom=0.95in]{geometry}
\newcommand{\tauCohomology}{T}
\newcommand{\FCohomology}{S}


\theoremstyle{theorem}
%\newtheorem*{mainthm*}{Main Theorem}

\newcommand{\note}[1]{\marginpar{\sffamily\tiny #1}}

\def\todo#1{\textcolor{Mahogany}%
{\sffamily\footnotesize\newline{\color{Mahogany}\fbox{\parbox{\textwidth-15pt}{#1}}}\newline}}

\renewcommand{\O}{\mathcal O}
\newcommand{\ie}{{\itshape i.e.} }
\newcommand{\roundup}[1]{\lceil #1 \rceil}
\newcommand{\rounddown}[1]{\lfloor #1 \rfloor}
\renewcommand{\to}[1][]{\xrightarrow{\ #1\ }}
\newcommand{\onto}[1][]{\protect{\xrightarrow{\ #1\ }\hspace{-0.8em}\rightarrow}}
\newcommand{\into}[1][]{\lhook \joinrel \xrightarrow{\ #1\ }}
%\newcommand{DBDual}[1]{{\underline{\omega}_{#1}}}
\newcommand{\DBDual}[1]{{{\uuline \omega}{}^{\mydot}_{#1}}}
\newcommand{\Tt}{{\mathfrak{T}}}
%\newcommand{\bn}{{\mathfrak{n}} }
\newcommand{\sol}[1]{ \vskip 6pt{\color{blue} {\bf Solution:} #1} \vskip 6pt}

\newcommand{\bR}{{\mathbb{R}}}
\newcommand{\va}{{\vec{a}}}
\newcommand{\vb}{{\vec{b}}}
\newcommand{\vzero}{{\vec{0}}}
\newcommand{\vi}{{\vec{i}}}
\newcommand{\vj}{{\vec{j}}}
\newcommand{\vk}{{\vec{k}}}
\newcommand{\vh}{{\vec{h}}}
\newcommand{\vr}{{\vec{r}}}
\newcommand{\vu}{{\vec{u}}}
\newcommand{\vv}{{\vec{v}}}
\newcommand{\vw}{{\vec{w}}}
\newcommand{\vx}{{\vec{x}}}
\newcommand{\vy}{{\vec{y}}}
\newcommand{\vz}{{\vec{z}}}
\newcommand{\vT}{{\vec{T}}}
\newcommand{\vN}{{\vec{N}}}
\newcommand{\vF}{{\vec{F}}}
\newcommand{\vG}{{\vec{G}}}
\newcommand{\bF}{\mathbb{F}}
\newcommand{\bQ}{\mathbb{Q}}
\newcommand{\bZ}{\mathbb{Z}}
\newcommand{\bC}{\mathbb{C}}


\begin{document}
\title{Worksheet \#2-- Math 5320\\Fall 2019}
\author{Due Friday, January 24th, in Gradescope}

\maketitle
You may work in groups of up to 3 on this assignment.  Only one assignment needs to be turned in per group.  It still needs to be turned in on gradescope.  Many of these problems are also problems in the textbook.

\noindent
{\bf 1.}  Find generators for the kernels of the following ring homomorphisms.  In other words, write each kernel as an ideal of the form $\langle f, g, \dots\rangle$ where $f, g, \dots$ are elements of the domain.
\vskip 3pt
{\bf (a)}  $\bR[x,y] \to \bR$, defined by $f(x,y) \mapsto f(0,0)$.
\vskip 110pt
{\bf (b)}  $\bR[x] \to \bC$, defined by $f(x) \mapsto f(2+i)$.
\vskip 110pt
{\bf (c)}  $\bZ[x] \to \bR$, defined by $f(x) \mapsto f(1+\sqrt{2})$.
\vskip 110pt
{\bf (d)}  $\bC[x,y,z] \to \bC[t]$ defined by $x \mapsto t, y \mapsto t^2, z \mapsto t^3$.
\newpage
\noindent
{\bf 2.}  Suppose we have two ideals $I, J \subseteq R$.  Then consider the sets
\begin{itemize}
    \item[(a)] $I + J = \{ x + y \;|\; x \in I, y \in J\}$.
    \item[(b)] $I \cap J$.
    \item[(c)] $I \cup J$.
    \item[(d)] $\{x\cdot y \;|\; x \in I, y \in J\}$.
    \item[(e)] $I \cdot J = \big\{\sum_{\text{finite}} x_i \cdot y_i \;\big|\; x_i \in I, y_i \in J \big\}$.
\end{itemize} 
Determine which are also ideals.  Prove it or give a counter example.
\newpage
\noindent
{\bf 3.}  Two ideals $I, J \subseteq R$ are called \emph{strongly coprime} if $I + J = R$.  Show that if $I$ and $J$ are strongly coprime, then $I \cap J = I \cdot J$.  Give an example to show that this conclusion is false without the strongly coprime hypothesis.  
\newpage
\noindent
{\bf 4.}  Suppose $I,J \subseteq R$ are two strongly coprime ideals.  Prove that for any $a,b \in R$ the system of congruences
\[
    \begin{array}{rcl}
        x + I & = & a + I\\
        x + J & = & b + J
    \end{array}
\]
has a solution and that if $x, x' \in R$ are two solutions, then $x + I \cap J = x' + I \cap J$.
\newpage
\noindent
{\bf 5.}  Suppose that $I, J \subseteq R$ are ideals.  Let $\phi : R \to (R/I) \times (R/J)$ be the map that sends $r \in R$ to the tuple $(r + I, r+J)$.  Show that the kernel of $\phi$ is $I \cap J$.  Then use the first isomorphism theorem and the previous problem to conclude that if $I$ and $J$ are coprime, that $R /(I \cap J) \cong (R/I) \times (R/J)$.
\newpage
\noindent
{\bf 6.}  Suppose that $I, J, K \subseteq R$ are ideals that are pairwise strongly coprime.  Show that $I + (J \cap K) = R$ as well.  Conclude by appropriately stating and generalizing problem {\bf 4.} to systems of 3 (or more) equations.
\end{document} 